The canonical, Galois-equivariant identification
\cite[Part I, Theorem 12.1(a)]{milne-av}
\[
 V_\ell(A)\simeq
 H^1_{\mathrm{et}}(A_{\overline F},\mathbb Q_\ell)^\vee
\]
gives the same property for the Tate module. De Rham representations
are closed under duals and subquotients
\cite[Theorem 5.2.1 and Proposition 6.3.3]{brinon-conrad}.
Extension of the coefficient field is, on underlying
$\mathbb Q_\ell$-representations, a finite direct sum.
The displayed decomposition therefore makes $r|_{G_F}$ de Rham.

Here finite-extension descent requires no semistable descent
assertion. To see it explicitly, pass to a finite Galois extension
$F'/K$ containing $F$, and let $W$ be the underlying
$\mathbb Q_\ell$-representation of $r$. The vector space
\[
 D_{F'}=(B_{\mathrm{dR}}\otimes_{\mathbb Q_\ell}W)^{G_{F'}}
\]
has its natural semilinear $\mathrm{Gal}(F'/K)$-action, and its
invariants are
$D_K=(B_{\mathrm{dR}}\otimes_{\mathbb Q_\ell}W)^{G_K}$.
Finite Galois descent gives $F'\otimes_K D_K\simeq D_{F'}$,
also on each filtration step. Since $W|_{G_{F'}}$ is de Rham,
$\dim_K D_K=\dim_{F'}D_{F'}=\dim_{\mathbb Q_\ell}W$.
This proves the assertion over $K$; it is the finite-extension
case of \cite[Proposition 6.3.8]{brinon-conrad}.

The two Hodge numbers of $H^1_{\mathrm{dR}}(A/F)$ are $2,2$.
Consequently the Tate module has weights $0,0,-1,-1$.
At each coefficient embedding, its two isomorphic summands
have the same weight multiset, so each summand has weights $0,-1$.
The filtration-compatible descent preserves these weights.
\end{proof}
